\documentclass[11pt,a4paper]{article}
\usepackage[T1]{fontenc}
\usepackage[utf8]{inputenc}
\usepackage{amsmath,amssymb}
\usepackage{mathpazo}
\usepackage{courier}
\usepackage{babel}
\babelprovide[import,main]{ngerman}
\usepackage[a4paper,margin=2.5cm]{geometry}
\usepackage{booktabs}
\usepackage{enumitem}
\usepackage{microtype}
\usepackage{xcolor}
\usepackage[hidelinks]{hyperref}
\setlength{\parindent}{0pt}
\setlength{\parskip}{0.55em}
\setlength{\emergencystretch}{3em}
\renewcommand{\contentsname}{Inhalt}
\renewcommand{\tablename}{Tabelle}
\newcommand{\re}{\operatorname{Re}}
\newcommand{\im}{\operatorname{Im}}
\newcommand{\E}{\operatorname{E}}
\newcommand{\Qd}{Q^{\mathrm{DAC}}}
\newcommand{\Qa}{Q^{\mathrm{ADC}}}
\newcommand{\file}[1]{\texttt{#1}}
\newcommand{\schritt}[1]{\subsection*{Schritt #1}}
\definecolor{note}{gray}{0.35}

\title{\textbf{Digitales Beamforming mit DAC-Quantisierung}\\[0.3em]
\large Berechnung der MI-Kurven, Schritt für Schritt}
\author{Gearbox-PHY, Studie analoges gegen digitales Beamforming}
\date{12. Oktober 2026, ergänzt um den digitalen LOS-Empfänger und die Präkodierung}

\begin{document}
\maketitle

\noindent Dieses Dokument beschreibt, wie die Kurven der Transinformation (MI) für digitales
Eigen-Beamforming mit endlicher DAC-Auflösung entstehen. Es folgt dem Code in
\file{QuantizedMimoMI}:
\begin{itemize}[leftmargin=1.4em,itemsep=0.05em,topsep=0.2em]
\item \file{dacQuantize.m} und \file{ergodicMiBeamforming.m} (Quantisierer und MI),
\item \file{runQamSweepBfRayleighDac.m} (Rayleigh),
\item \file{runQamSweepBfLosDigitalDac.m} und \file{runQamSweepBfIdealDac.m} (LOS).
\end{itemize}
Abschnitt~\ref{sec:prec} erklärt, wie die Gewichte gefunden werden, Abschnitt~\ref{sec:ray}
behandelt den Rayleigh-Kanal, Abschnitt~\ref{sec:los} die LOS-Referenz mit beiden
Empfängermodellen, Abschnitt~\ref{sec:gb} die Übergabe an das Gearbox,
Abschnitt~\ref{sec:annahmen} die Annahmen.

\section{Größen und Bezeichnungen}

\begin{tabular}{@{}ll@{}}
$N_t,\,N_r$ & Zahl der Sende- und Empfangsantennen (hier $N_t=N_r=N$)\\
$\mathcal X$ & quadratisches $M$-QAM-Alphabet, gleich wahrscheinlich, $\E|x|^2=1$\\
$H\in\mathbb C^{N_r\times N_t}$ & Kanalmatrix, dem Sender und dem Empfänger bekannt\\
$\rho$ & SNR bezogen auf die \emph{gesamte} Sendeleistung; Rauschvarianz $\sigma^2=1/\rho$\\
$b_{\mathrm{DAC}}=\tfrac12\log_2M+k$ & DAC-Bit je I/Q-Zweig und Antenne, Stufe $k\in\{0,1,2,3\}$\\
$L=2^{b_{\mathrm{DAC}}}$ & Stufenzahl des DAC je Zweig\\
$B=\tfrac12\log_2M+3$ & ADC-Bit je I/Q-Zweig und Antenne (Regel \file{fixedB}); $L_A=2^B$\\
\end{tabular}

Ohne DAC-Quantisierung ($k$ beliebig groß) ist das Verfahren das bisherige: der Sender strahlt
$t(x)=v\,x$ ab. Mit $b_{\mathrm{DAC}}=\tfrac12\log_2M$ Bit lässt sich ein \emph{ungedrehtes}
QAM-Symbol exakt darstellen, das Produkt $v_i\,x$ mit komplexem Gewicht $v_i$ im Allgemeinen nicht.

\section{Präkodierung: wie die Gewichte entstehen}\label{sec:prec}

\subsection{Aufgabe}

Gesendet wird \emph{ein} Datenstrom. Jede Antenne strahlt dasselbe Symbol $x$ ab, multipliziert
mit einem eigenen komplexen Gewicht:
\begin{equation}
t=v\,x,\qquad v\in\mathbb C^{N_t},\qquad \|v\|=1 .
\end{equation}
Die Nebenbedingung $\|v\|=1$ hält die Gesamtleistung fest: $\E\|t\|^2=\|v\|^2\,\E|x|^2=1$. Betrag
und Phase jedes $v_i$ sind frei wählbar. Genau das unterscheidet die digitale Präkodierung von der
analogen, bei der nur die Phasen frei sind ($|v_i|=1/\sqrt{N_t}$), und es ist der Grund, warum jede
Antenne ihren eigenen DAC braucht.

Am Empfänger kommt ohne Quantisierung
\begin{equation}
y=H\,v\,x+n=h_{\mathrm{eff}}\,x+n,\qquad h_{\mathrm{eff}}=Hv\in\mathbb C^{N_r}
\end{equation}
an. Ein Empfänger, der die $N_r$ Signale mit $w=h_{\mathrm{eff}}/\|h_{\mathrm{eff}}\|$ kombiniert
(Maximum-Ratio-Combining), erhält den skalaren Kanal
\begin{equation}
z=w^{\mathsf H}y=\|Hv\|\,x+\tilde n,\qquad \tilde n\sim\mathcal{CN}(0,\sigma^2),\qquad
\mathrm{SNR}_{\mathrm{eff}}=\rho\,\|Hv\|^2 .
\end{equation}
Die Wahl der Gewichte ist damit ein Optimierungsproblem: der Arraygewinn $\|Hv\|^2$ soll so groß
wie möglich werden,
\begin{equation}\label{eq:opt}
v^\star=\arg\max_{\|v\|=1}\;\|Hv\|^2=\arg\max_{\|v\|=1}\;v^{\mathsf H}H^{\mathsf H}H\,v .
\end{equation}

\subsection{Lösung über die Singulärwertzerlegung}

Der Kanal wird zerlegt in
\begin{equation}
H=U\Sigma V^{\mathsf H}=\sum_{j=1}^{r}\sigma_j\,u_jv_j^{\mathsf H},\qquad
\sigma_1\ge\sigma_2\ge\dots\ge\sigma_r>0,
\end{equation}
mit orthonormalen Senderichtungen $v_j\in\mathbb C^{N_t}$, orthonormalen Empfangsrichtungen
$u_j\in\mathbb C^{N_r}$ und dem Rang $r$. Jedes Paar $(v_j,u_j)$ ist ein \emph{Eigenmode}: was in
Richtung $v_j$ gesendet wird, kommt in Richtung $u_j$ an, verstärkt um $\sigma_j$.

Jeder Einheitsvektor lässt sich als $v=\sum_jc_jv_j+v_\perp$ schreiben, mit
$\sum_j|c_j|^2\le1$ und einem Anteil $v_\perp$ im Nullraum von $H$. Dann ist
\begin{equation}
\|Hv\|^2=\Bigl\|\sum_j\sigma_jc_j\,u_j\Bigr\|^2=\sum_j\sigma_j^2\,|c_j|^2\;\le\;\sigma_1^2 ,
\end{equation}
mit Gleichheit genau dann, wenn die ganze Leistung im stärksten Eigenmode liegt
($|c_1|=1$). Die Lösung von \eqref{eq:opt} ist also der \emph{dominante rechte Singulärvektor},
\begin{equation}
v^\star=v_1,\qquad \|Hv_1\|^2=\sigma_1^2=\lambda_{\max}\bigl(H^{\mathsf H}H\bigr),\qquad Hv_1=\sigma_1u_1 .
\end{equation}
Gleichwertig: $v_1$ ist der Eigenvektor von $H^{\mathsf H}H$ zum größten Eigenwert. Daher der Name
Eigen-Beamforming. Das Signal kommt entlang $u_1$ an; der passende Kombinierer ist $w=u_1$.

Im Code ist das eine Zeile je Kanalrealisierung: die Zerlegung \file{[\textasciitilde,\textasciitilde,V] = svd(H,'econ')}
und daraus die erste Spalte von $V$. Für $N_t=1$ gibt es nichts zu wählen, $v=1$.

\subsection{Was die Gewichte anschaulich tun}

\textbf{Eine Empfangsantenne ($N_r=1$).} Mit $H=h^{\mathsf T}$ ist
\begin{equation}
v_i=\frac{h_i^\ast}{\|h\|},\qquad \|Hv\|^2=\|h\|^2=\sum_i|h_i|^2 .
\end{equation}
Die Phase jedes Gewichts hebt die Phase seines Pfades auf, sodass alle Beiträge am Empfänger
gleichphasig ankommen. Der Betrag ist proportional zu $|h_i|$: ein starker Pfad bekommt mehr
Leistung als ein schwacher. Das ist Maximum-Ratio-Transmission.

\textbf{LOS, Rang 1.} Für $H=a_ra_t^{\mathsf H}$ gibt es nur einen Eigenmode,
\begin{equation}
v_1=\frac{a_t}{\sqrt{N_t}},\qquad u_1=\frac{a_r}{\sqrt{N_r}},\qquad \sigma_1^2=N_tN_r .
\end{equation}
Alle Gewichte haben denselben Betrag und unterscheiden sich nur in der Phase
$e^{\mathrm j\pi(i-1)u_t}$: die Antennenzeile wird in die Richtung $u_t$ ausgerichtet. Der Gewinn ist
der volle $N_tN_r$. Hier fällt die digitale mit der analogen Lösung zusammen.

\textbf{Rayleigh, voller Rang.} Die Einträge von $H$ sind unabhängig, es gibt $N$ Eigenmoden mit
verschiedenen $\sigma_j$. Die Gewichte haben verschiedene Beträge \emph{und} Phasen; ihre Beträge
streuen von Antenne zu Antenne. Der mittlere Gewinn $\E[\sigma_1^2]$ liegt unter $N_tN_r$, weil
sich die Kanalenergie auf alle Moden verteilt: um etwa $0{,}6$~dB bei $2\times2$, $2{,}2$~dB bei
$4\times4$, $4{,}3$~dB bei $8\times8$ und $6{,}8$~dB bei $16\times16$. Die übrigen Moden bleiben beim
Beamforming ungenutzt; Multiplexing würde gerade sie für weitere Datenströme verwenden.

\subsection{Die freie Phase}

Mit $v$ löst auch $v\,e^{\mathrm j\theta}$ das Problem \eqref{eq:opt}, für jedes $\theta$: der
Gewinn ändert sich nicht, am Empfänger ist nur die Konstellation um $\theta$ gedreht. Ohne
Quantisierung ist $\theta$ deshalb bedeutungslos. Mit Quantisierung nicht mehr, denn $\theta$
legt fest, wie die Symbole auf den quadratischen Rastern von DAC und ADC liegen. Der Code nimmt
eine feste Konvention,
\begin{equation}
v=v_1\,e^{-\mathrm j\arg([v_1]_1)},
\end{equation}
sodass das erste Gewicht reell und positiv ist. Antenne~1 sendet damit ein ungedrehtes Symbol.
Die Phase wird nicht auf die MI optimiert; eine Auswahl der besten aus 16 Phasen brächte am
Quantisierer höchstens etwa $0{,}2$~dB, und nur bei $k=0$.

\subsection{Was dabei vorausgesetzt und was nicht optimiert ist}

\begin{itemize}[leftmargin=1.4em,itemsep=0.2em,topsep=0.2em]
\item \textbf{Der Sender kennt $H$ vollständig und fehlerfrei.} Kanalschätzung und Rückmeldung
kosten im Modell nichts.
\item \textbf{$v_1$ ist optimal für einen Strom ohne Quantisierung.} Dort hängt die Rate nur von
$\|Hv\|^2$ ab. Mit ADCs und endlichem Alphabet ist das nicht mehr streng richtig; $v_1$ bleibt die
naheliegende Wahl, ist aber nicht als MI-optimal bewiesen.
\item \textbf{Erst gewichten, dann runden.} Die Gewichte werden so bestimmt, als gäbe es keinen DAC;
die Rundung kommt danach (Schritt~4). Eine Präkodierung, die gleich die auf dem Raster
darstellbaren Sendevektoren sucht, würde den DAC von vornherein berücksichtigen und wäre bei
wenigen Bit besser.
\item \textbf{In der MI steht kein Kombinierer.} Berechnet wird $I(x;y_q\mid H)$ aus allen
quantisierten Empfangssignalen; das ist die Rate des besten digitalen Empfängers. Der Kombinierer
$w=u_1$ oben dient nur der Herleitung der Gewichte.
\item \textbf{Die Beträge der Gewichte bestimmen die Aussteuerung.} Weil $|v_i|$ im Rayleigh-Kanal
streut, hat jede Antenne einen anderen Spitzenwert; deshalb wird der DAC je Antenne ausgesteuert
(Schritte~3 und~4).
\end{itemize}

\section{Rayleigh-Kanal}\label{sec:ray}

Für jede Kurve $(N,M,k)$ und jeden SNR-Punkt $\rho$ werden $K=400$ Kanäle $H_1,\dots,H_K$ mit
unabhängigen Einträgen $\mathcal{CN}(0,1)$ gezogen (Threefry-Strom $\kappa$ für $H_\kappa$, Seed 42 -- dieselben
Kanäle wie in den Kurven ohne DAC). Die Schritte 1 bis 7 gelten je Kanal.

\schritt{1: Gewichte}
Wie in Abschnitt~\ref{sec:prec} hergeleitet: Singulärwertzerlegung $H=U\Sigma V^{\mathsf H}$ mit
$\sigma_1\ge\sigma_2\ge\dots$, dominanter rechter Singulärvektor $v_1$. Die freie gemeinsame Phase
wird durch die feste Konvention entfernt:
\begin{equation}
v \;=\; v_1\,e^{-\mathrm j\arg([v_1]_1)},\qquad \|v\|=1,\qquad [v]_1\in\mathbb R_{>0}.
\end{equation}
Antenne~1 sendet damit immer ein ungedrehtes Symbol. Die Phase ist nicht optimiert.

\schritt{2: Sollsignal je Antenne}
\begin{equation}
u_i(x)\;=\;v_i\,x,\qquad i=1,\dots,N_t,\quad x\in\mathcal X .
\end{equation}
Je Antenne sind das $M$ komplexe Sollwerte: die um $\arg v_i$ gedrehte und mit $|v_i|$ skalierte
Konstellation.

\schritt{3: Spitzenwert}
\begin{equation}
A_i\;=\;\max_{x\in\mathcal X}\,\max\bigl(|\re u_i(x)|,\;|\im u_i(x)|\bigr).
\end{equation}

\schritt{4: Runden auf das DAC-Raster}
Für eine Stufenzahl $L_c$ und einen Spitzenwert $A$ ist das Raster
\begin{equation}
\mathcal G(A,L_c)=\Bigl\{\,A\,\frac{2m-(L_c-1)}{L_c-1}\;:\;m=0,\dots,L_c-1\Bigr\},
\end{equation}
also $L_c$ gleichabständige Stufen, die äußersten bei $\pm A$. Gerundet wird auf die nächste Stufe,
Real- und Imaginärteil getrennt:
\begin{align}
m^\ast(r)&=\min\Bigl\{L_c-1,\;\max\Bigl\{0,\;\operatorname{round}\Bigl(\frac{r\,(L_c-1)/A+(L_c-1)}{2}\Bigr)\Bigr\}\Bigr\},\\
q_{A,L_c}(r)&=A\,\frac{2m^\ast(r)-(L_c-1)}{L_c-1},\qquad
\Qd_{A,L_c}(z)=q_{A,L_c}(\re z)+\mathrm j\,q_{A,L_c}(\im z).
\end{align}
Wegen $|r|\le A$ wird nie abgeschnitten.

\textbf{Wählbare Aussteuerung.} Zugelassen sind alle geraden Stufenzahlen
$L_c\in\{2,4,\dots,L\}$. Je Antenne wird diejenige mit dem kleinsten quadratischen Fehler genommen,
bei Gleichstand die größte:
\begin{equation}
e_i(L_c)=\sum_{x\in\mathcal X}\bigl|\Qd_{A_i,L_c}(u_i(x))-u_i(x)\bigr|^2,\qquad
L_i^\ast=\arg\min_{L_c\in\{2,4,\dots,L\}}e_i(L_c),\qquad
\tilde t_i(x)=\Qd_{A_i,L_i^\ast}\bigl(u_i(x)\bigr).
\end{equation}

\emph{Warum das ein $L$-Stufen-DAC kann.} Ein DAC mit $L$ Stufen und Vollausschlag
$F=A\,(L-1)/(L_c-1)$ hat die Stufen
\[
F\,\frac{2n-(L-1)}{L-1}=A\,\frac{2n-(L-1)}{L_c-1},\qquad n=0,\dots,L-1 .
\]
Der Zähler $2n-(L-1)$ durchläuft alle ungeraden Zahlen von $-(L-1)$ bis $L-1$ ($L$ ist gerade). Für
gerades $L_c\le L$ ist $2m-(L_c-1)$ ebenfalls ungerade und betragsmäßig höchstens $L_c-1\le L-1$.
Die Menge $\mathcal G(A,L_c)$ besteht also genau aus den inneren $L_c$ Stufen dieses DAC: ein
feinerer DAC bildet jeden gröberen mit gerader Stufenzahl exakt nach, indem er seinen
Vollausschlag vergrößert.

\emph{Folgen.}
\begin{itemize}[leftmargin=1.4em,itemsep=0.15em,topsep=0.2em]
\item Der Fehler je Antenne wächst nie mit der Stufe $k$, denn die Kandidatenmenge wird nur größer.
\item Ein ungedrehtes QAM-Symbol ($v_i$ reell) ist für jedes $k\ge0$ exakt: $L_c=\sqrt M$ ergibt
$e_i=0$. Mit festem $L_c=L$ wäre das nicht so (16-QAM ungedreht, 3~Bit: 8{,}3\,\% EVM).
\item Für $N_t=1$ ist $v=1$, nichts ist gedreht; der Code überspringt den Quantisierer, die Kurve
ist bitgleich zur bisherigen, und der DAC kostet $\tfrac12\log_2M$ Bit.
\end{itemize}

\emph{Vergleichsfall \quotedblbase gemeinsame Aussteuerung\textquotedblleft.} Ein Spitzenwert $A=\max_iA_i$ und eine
Stufenzahl für alle Antennen. Gewählt wird nach dem Gesamtfehler $\sum_ie_i(L_c)$. Er kommt ohne
Verstärkungseinstellung je Antenne hinter dem DAC aus.

\schritt{5: Normierung der Sendeleistung}
Mit $\tilde t(x)=[\tilde t_1(x),\dots,\tilde t_{N_t}(x)]^{\mathsf T}$:
\begin{equation}
\alpha=\Bigl(\frac1M\sum_{x\in\mathcal X}\|\tilde t(x)\|^2\Bigr)^{-1/2},\qquad
t(x)=\alpha\,\tilde t(x),\qquad \frac1M\sum_{x\in\mathcal X}\|t(x)\|^2=1 .
\end{equation}
Die mittlere Gesamtleistung ist damit dieselbe wie ohne DAC, und $\rho$ bleibt die Gesamt-SNR.
Zu jedem der $M$ Symbole gehört jetzt ein fester Sendevektor $t(x)$.

\schritt{6: Kanal und ADC}
\begin{equation}
y=H\,t(x)+n,\qquad n\sim\mathcal{CN}(0,\sigma^2I_{N_r}),\qquad \mu(x)=H\,t(x).
\end{equation}
Jede Empfangsantenne hat einen gleichförmigen $B$-Bit-ADC je I/Q-Zweig. Sein Bereich folgt der
Empfangsleistung bei gegebenem $H$:
\begin{equation}
\sigma_y^2=\frac{1}{2MN_r}\sum_{x\in\mathcal X}\|\mu(x)\|^2+\frac{\sigma^2}{2},\qquad V=3\,\sigma_y ,
\end{equation}
mit den Schwellen $\tau_0=-\infty$, $\tau_\ell=-V+2V\ell/L_A$ für $\ell=1,\dots,L_A-1$ und
$\tau_{L_A}=+\infty$. Beobachtet wird $y_q$: für jede Antenne und jeden Zweig der Index der Zelle
$[\tau_\ell,\tau_{\ell+1})$, in der der Wert liegt. Die ADC-Seite ist gegenüber den Kurven ohne DAC
unverändert.

\schritt{7: Transinformation bei gegebenem Kanal}
\begin{equation}
I(x;y_q\mid H)=\mathsf H(y_q\mid H)-\mathsf H(y_q\mid x,H).
\end{equation}
Für eine reelle Komponente mit Mittelwert $m$ und Rausch-Standardabweichung
$\sigma_r=\sigma/\sqrt2$ ist die Wahrscheinlichkeit der Zelle $\ell$
\begin{equation}
p_\ell(m)=\Phi\Bigl(\frac{\tau_{\ell+1}-m}{\sigma_r}\Bigr)-\Phi\Bigl(\frac{\tau_\ell-m}{\sigma_r}\Bigr),
\end{equation}
mit der Verteilungsfunktion $\Phi$ der Standardnormalverteilung. Gegeben $x$ sind die $2N_r$
reellen Komponenten unabhängig:
\begin{equation}
P(y_q\mid x)=\prod_{i=1}^{N_r}p_{\ell_i^{\mathrm R}}\bigl(\re\mu_i(x)\bigr)\,p_{\ell_i^{\mathrm I}}\bigl(\im\mu_i(x)\bigr),\qquad
P(y_q)=\frac1M\sum_{x\in\mathcal X}P(y_q\mid x).
\end{equation}
Die bedingte Entropie ist damit geschlossen berechenbar,
\begin{equation}
\mathsf H(y_q\mid x,H)=\frac1M\sum_{x\in\mathcal X}\sum_{i=1}^{N_r}\Bigl[h\bigl(p(\re\mu_i(x))\bigr)+h\bigl(p(\im\mu_i(x))\bigr)\Bigr],\qquad
h(p)=-\sum_\ell p_\ell\log_2p_\ell .
\end{equation}
Für $\mathsf H(y_q\mid H)$ gibt es zwei Wege, je nach Größe des Ausgaberaums:
\begin{itemize}[leftmargin=1.4em,itemsep=0.15em,topsep=0.2em]
\item \textbf{exakt}, falls $M\cdot L_A^{2N_r}\le5\cdot10^7$: Summe
$-\sum_{y_q}P(y_q)\log_2P(y_q)$ über alle Ausgaben;
\item \textbf{quasi-exakt} sonst: $S=2000$ Beobachtungen $y_q^{(s)}$ werden gezogen
($x$ gleichverteilt, $n$ wie oben), und
$\mathsf H(y_q\mid H)\approx-\frac1S\sum_{s=1}^S\log_2P\bigl(y_q^{(s)}\bigr)$, wobei $P(\cdot)$ als
Summe über \emph{alle} $M$ Kandidaten exakt ausgewertet wird.
\end{itemize}
Der einzige Unterschied zum Fall ohne DAC ist die Kandidatenmenge: $\{H\,t(x)\}$ statt
$\{H v\,x\}$. Die Rundung zählt als verändertes Alphabet, nicht als Rauschen.

\schritt{8: Mittelung über den Kanal}
\begin{equation}
I(\rho)=\frac1K\sum_{\kappa=1}^K I(x;y_q\mid H_\kappa),\qquad K=400 ,
\end{equation}
mit dem 95-\%-Intervall $1{,}96\,s/\sqrt K$ aus der Streuung $s$ über die Kanäle. Das Raster ist
$\rho\in\{-15,-14,\dots,25\}$~dB.

\section{LOS-Referenz}\label{sec:los}

Der Kanal hat Rang~1,
\begin{equation}
H=a_r\,a_t^{\mathsf H},\qquad a(u)_n=e^{\mathrm j\pi nu},\quad n=0,\dots,N-1,\quad u=\sin\theta ,
\end{equation}
für eine gleichförmige Antennenzeile mit halber Wellenlänge Abstand, mit der Senderichtung $u_t$
in $a_t$ und der Empfangsrichtung $u_r$ in $a_r$. Die Gewichte sind reine Phasen,
$v=a_t/\sqrt{N_t}$; das erste Element ist reell. Die Schritte 2 bis 5 laufen wie im
Rayleigh-Fall und liefern $t(x)$. Für den Empfänger gibt es zwei Modelle.

\subsection{Digitaler Empfänger: ein ADC je Antenne}\label{sec:losd}

Rein digitales Beamforming hat $N_r$ Empfangsketten und damit $N_r$ ADCs; kombiniert wird
\emph{nach} der Quantisierung. Das ist derselbe Empfänger wie im Rayleigh-Fall, und gerechnet wird
mit denselben Schritten 6 und 7, nur mit dem deterministischen Kanal oben:
\begin{equation}
y=a_r\,s(x)+n,\qquad s(x)=a_t^{\mathsf H}\,t(x),\qquad \mu_i(x)=e^{\mathrm j\pi(i-1)u_r}\,s(x),
\end{equation}
mit einem $B$-Bit-ADC je Antenne und $I(x;y_q\mid H)$ wie in Schritt~7. Alle Antennen sehen
dasselbe skalare Signal $s(x)$, jede um ihre Phase gedreht und mit eigenem Rauschen. Die
Empfangsrichtung zählt hier, weil sie bestimmt, wie die Signale der Antennen auf den ADC-Rastern
liegen. Ohne DAC-Quantisierung ist $s(x)=\sqrt{N_t}\,x$.

Gemittelt wird über $32$ Richtungspaare,
\begin{equation}
I(\rho)=\frac1{32}\sum_{j=1}^{32}I\bigl(u_{t,j},u_{r,j}\bigr),\qquad
u_{t,j}=\frac{j-\tfrac12}{32},\qquad u_{r,j}=u_{t,\pi(j)},\quad \pi(j)=\bigl(13\,(j-1)\bmod 32\bigr)+1 .
\end{equation}
Die Empfangsrichtungen sind dieselben Werte in anderer Reihenfolge, damit Sende- und
Empfangsrichtung nicht gekoppelt sind. Die MI ändert sich nicht, wenn beide Vorzeichen zugleich
wechseln, deshalb genügt $(0,1)$. Der Kanal ist deterministisch; im quasi-exakten Fall werden je
Richtungspaar vier Auswertungen gemittelt, um das Rauschen des Schätzers zu senken. Gerechnet
werden die Kurve ohne DAC-Quantisierung und die Stufen $k=0,\dots,3$
(\file{runQamSweepBfLosDigitalDac.m}). Das SNR-Raster ist Element-SNR $-20$ bis $45$~dB in
Schritten von $1$~dB, Gesamt-SNR $=$ Element-SNR $-10\log_{10}(N_tN_r)$.

\subsection{Vereinfachung: Kombinieren vor einem ADC}\label{sec:los1}

Die bisherige LOS-Referenz kombiniert die $N_r$ Signale kohärent, \emph{bevor} ein einziger ADC
quantisiert. Das ist der \emph{analoge} Empfänger. Für die digitale Seite ist es eine
Vereinfachung; exakt beschreibt es den Fall \quotedblbase  Sender digital, Empfänger analog\textquotedblleft.
Übrig bleibt der skalare Kanal
\begin{equation}
r=c(x)+n,\qquad c(x)=\sqrt{N_r}\;a_t^{\mathsf H}\,t(x),\qquad n\sim\mathcal{CN}(0,\sigma^2).
\end{equation}
Ohne DAC-Quantisierung ist $c(x)=\sqrt{N_tN_r}\,x$: die SISO-Kurve mit Gewinn $N_tN_r$. Mit
Quantisierung ist $\{c(x)\}$ ein verzerrtes Alphabet, dessen I- und Q-Komponente nicht mehr
unabhängig sind. Die MI ist trotzdem exakt berechenbar: mit den Zeilenvektoren
$p^{\mathrm R}(x)=p(\re c(x))$ und $p^{\mathrm I}(x)=p(\im c(x))$ der Länge $L_A$ gilt
\begin{equation}
P(\ell^{\mathrm R},\ell^{\mathrm I})=\frac1M\sum_{x\in\mathcal X}p^{\mathrm R}_{\ell^{\mathrm R}}(x)\,p^{\mathrm I}_{\ell^{\mathrm I}}(x),\qquad
I(u_t)=\mathsf H(P)-\frac1M\sum_{x\in\mathcal X}\bigl[h(p^{\mathrm R}(x))+h(p^{\mathrm I}(x))\bigr].
\end{equation}
Die Schwellen folgen aus $\sigma_y^2=\tfrac{1}{2M}\sum_x|c(x)|^2+\tfrac{\sigma^2}{2}$ wie in Schritt~6.
Gemittelt wird über die $32$ Senderichtungen $u_{t,j}$; die Empfangsrichtung fällt durch das
Kombinieren heraus. Für $u_t=0$ sind alle Gewichte gleich und reell; nichts ist gedreht, und die
Kurve ist die bisherige (\file{runQamSweepBfIdealDac.m}, Element-SNR in Schritten von $0{,}5$~dB).

\subsection{Unterschied der beiden Empfänger}

Mit $N_r$ ADCs mitteln sich die Quantisierungsfehler der Antennen beim digitalen Kombinieren
teilweise heraus; die Kurve liegt deshalb etwas \emph{über} der mit einem ADC. Aus den gerechneten
Kurven ohne DAC-Quantisierung (Stand 12.~Oktober 2026, $B=\tfrac12\log_2M+3$): der Vorteil bei
$I=0{,}75\log_2M$ beträgt $0{,}03$ bis $0{,}06$~dB für $2\times2$ und wächst auf $0{,}12$ bis
$0{,}17$~dB für $16\times16$. Für $2\times2$ mit 4-QAM und 16-QAM ist das exakt gerechnet
(höchstens $0{,}009$ bzw.\ $0{,}018$~bit); für größere Arrays ist der quasi-exakte Schätzer im
Spiel, dessen Rauschen hier bei rund $0{,}02$~bit liegt.

Der Mehrbedarf durch den DAC ist in beiden Modellen praktisch derselbe: die Werte der Tabelle in
Abschnitt~\ref{sec:erg} ändern sich um höchstens $0{,}02$~dB. Weitere Eigenmoden gibt es in beiden
Modellen nicht ($\sigma_j=0$ für $j\ge2$), und die Kurven mit DAC liegen auch mit $N_r$ ADCs nicht
über der ohne DAC (Überschreitungen im Rahmen des Schätzerrauschens).

\section{Zwei Sonderfälle zum Nachrechnen}

\textbf{4-QAM mit $k=0$.} Hier ist $L=2$; je Zweig bleibt nur das Vorzeichen. Mit
$v_i=|v_i|e^{\mathrm j\varphi_i}$ und $x=e^{\mathrm j(\pi/4+n\pi/2)}$ ist
\begin{equation}
\tilde t_i(x)=A_i\bigl(\operatorname{sgn}\re u_i(x)+\mathrm j\operatorname{sgn}\im u_i(x)\bigr)
=\sqrt2\,A_i\;e^{\mathrm j\frac\pi2\operatorname{round}(2\varphi_i/\pi)}\;x .
\end{equation}
Der DAC wirkt wie ein Phasenschieber mit 2~Bit: jede Antenne sendet weiter ein 4-QAM-Symbol, nur
mit einer auf Vielfache von $90^\circ$ gerundeten Phase. Das Empfangsalphabet bleibt exakt 4-QAM;
verloren geht Arraygewinn, nicht die Form der Konstellation.

\textbf{Viele Bit.} Für $k\to\infty$ geht $\tilde t_i(x)\to u_i(x)$, $\alpha\to1$ und $t(x)\to v\,x$:
die Kurve ohne DAC.

\section{Übergabe an das Gearbox}\label{sec:gb}

Jede Kurve wird mit ihrer DAC-Auflösung abgelegt (\file{sourceBdac}$\,=b_{\mathrm{DAC}}$, für
$N_t=1$ gleich $\tfrac12\log_2M$). Ein digitaler Sender bezahlt genau diese Auflösung,
\begin{equation}
P_{\mathrm{DAC}}=N_t\cdot2\Bigl[P_s\,(2^{b_{\mathrm{DAC}}}-1)+E_d\,b_{\mathrm{DAC}}\,B_{\mathrm w}\Bigr],
\end{equation}
mit der Bandbreite $B_{\mathrm w}$ und den Konstanten des gewählten DAC-Leistungsmodells
($P_s=15$, $5$ oder $1{,}2~\mu$W und $E_d=4{,}5$, $0{,}5$ oder $0{,}25$~pJ). Ein feinerer DAC
braucht weniger Sendeleistung und mehr DAC-Leistung. Für jeden Arbeitspunkt wird die günstigste
Stufe genommen:
\begin{equation}
E_{\mathrm{bit}}^\ast=\min_{k\in\{0,1,2,3\}}E_{\mathrm{bit}}^{(k)} .
\end{equation}
Ein analoger Sender hat einen einzigen DAC mit ungedrehtem Symbol; er bezahlt immer
$\tfrac12\log_2M$ Bit, die Drehung übernehmen die Phasenschieber.

\section{Annahmen}\label{sec:annahmen}

\begin{enumerate}[leftmargin=1.6em,itemsep=0.25em]
\item \textbf{Der Empfänger kennt $H$ und alle $M$ Sendevektoren $t(x)$.} Im Rayleigh-Fall und in
Abschnitt~\ref{sec:losd} wertet er alle $N_r$ quantisierten Signale gemeinsam aus; $I(x;y_q\mid H)$ ist die Rate des besten
digitalen Empfängers.
\item \textbf{Die Aussteuerung wird nach dem Fehler je Antenne gewählt}, nicht nach der MI.
\item \textbf{Kein Abschneiden.} Für ein gedrehtes QAM-Symbol brächte Abschneiden höchstens etwa
$0{,}1$~dB; bei mehreren präkodierten Strömen wäre es wichtig.
\item \textbf{Die gemeinsame Phase ist eine Konvention} (erstes Gewicht reell), nicht optimiert.
\item \textbf{Schlichtes Runden.} Eine Präkodierung, die die quantisierten Ausgänge direkt wählt,
wäre besser.
\item \textbf{Empfänger der LOS-Referenz.} Die digitale Seite gehört zu Abschnitt~\ref{sec:losd}
(ein ADC je Antenne). Die Vereinfachung aus Abschnitt~\ref{sec:los1} unterschätzt sie leicht.
\end{enumerate}

\textbf{Eine Folge von Annahme 1.} Mit $H=\sum_j\sigma_ju_jv_j^{\mathsf H}$ ist
\begin{equation}
H\,t(x)=\sigma_1u_1\,\bigl(v_1^{\mathsf H}t(x)\bigr)+\sum_{j\ge2}\sigma_ju_j\,\bigl(v_j^{\mathsf H}t(x)\bigr).
\end{equation}
Ohne DAC ist $t(x)=v\,x$ und die Summe verschwindet. Mit DAC liegt der Rundungsfehler auch in den
übrigen Eigenmoden und hängt von $x$ ab; ein Empfänger nach Annahme~1 gewinnt daraus Information.
In den Ergebnissen liegt deshalb im Rayleigh-Kanal die Kurve mit $k=0$ für $M\ge64$ und $N\ge4$
\emph{über} der Kurve ohne DAC (bis $0{,}10$~bit). Im LOS-Fall ist $\sigma_j=0$ für $j\ge2$, dort
tritt der Effekt nicht auf. Ein Empfänger, der nur auf $u_1$ kombiniert, hätte davon nichts.

\section{Ergebnis im Überblick}\label{sec:erg}

Mehrbedarf an SNR gegenüber der Kurve ohne DAC-Quantisierung bei $I=0{,}75\log_2M$, in dB,
Aussteuerung je Antenne (aus den gerechneten Kurven, Stand 9.~Oktober 2026). Die LOS-Zeilen
stammen aus der Vereinfachung mit einem ADC (Abschnitt~\ref{sec:los1}); mit einem ADC je
Antenne ergeben sich bis auf $0{,}02$~dB dieselben Werte:

\begin{center}
\begin{tabular}{llrrrrr}
\toprule
Kanal & $M$ & $N$ & $k=0$ & $k=1$ & $k=2$ & $k=3$\\
\midrule
LOS & 4 & 4 & 0{,}58 & 0{,}08 & 0{,}01 & 0{,}00\\
LOS & 4 & 16 & 0{,}83 & 0{,}11 & 0{,}01 & 0{,}00\\
LOS & 16 & 16 & 0{,}35 & 0{,}18 & 0{,}02 & 0{,}00\\
LOS & 256 & 16 & 0{,}02 & 0{,}01 & 0{,}02 & 0{,}00\\
\midrule
Rayleigh & 4 & 4 & 0{,}39 & 0{,}05 & 0{,}00 & 0{,}00\\
Rayleigh & 4 & 16 & 0{,}61 & 0{,}08 & 0{,}01 & 0{,}00\\
Rayleigh & 16 & 16 & $-0{,}06$ & 0{,}13 & 0{,}01 & 0{,}00\\
Rayleigh & 64 & 16 & $-0{,}27$ & 0{,}00 & 0{,}02 & 0{,}00\\
Rayleigh & 256 & 16 & $-0{,}29$ & $-0{,}05$ & $-0{,}01$ & 0{,}00\\
\bottomrule
\end{tabular}
\end{center}

Negative Werte sind die im vorigen Abschnitt beschriebene Folge von Annahme~1. Mit gemeinsamer
Aussteuerung beträgt der Mehrbedarf im Rayleigh-Kanal bei 4-QAM und $k=0$ zwischen $0{,}58$ und
$1{,}26$~dB.

{\color{note}\small Kontrollen im Code: \file{QuantizedMimoMI/qam/validate/validateBfDac.m} (elf
Tests, u.\,a. exakte Darstellung ungedrehter QAM, Normierung, $N_t=1$ bitgleich, 4-QAM mit $k=0$
bleibt 4-QAM, LOS bei $u=0$ gleich der bisherigen Kurve).}

\end{document}
